\documentclass[11pt,reqno]{amsart}
\usepackage[T1]{fontenc}
\usepackage{iftex}
\ifPDFTeX
  \usepackage[utf8]{inputenc}
\fi
\usepackage{lmodern}
\usepackage{amsmath,amssymb,amsthm,mathtools,mathrsfs}
\usepackage{booktabs}
\usepackage{enumitem}
\usepackage{microtype}
\usepackage{xcolor}
\usepackage[hidelinks]{hyperref}
\usepackage[nameinlink,noabbrev]{cleveref}

\numberwithin{equation}{section}
\raggedbottom

\newtheorem{theorem}{Theorem}[section]
\newtheorem{proposition}[theorem]{Proposition}
\newtheorem{lemma}[theorem]{Lemma}
\newtheorem{corollary}[theorem]{Corollary}
\newtheorem{question}[theorem]{Question}
\theoremstyle{definition}
\newtheorem{definition}[theorem]{Definition}
\newtheorem{example}[theorem]{Example}
\theoremstyle{remark}
\newtheorem{remark}[theorem]{Remark}

\newcommand{\A}{\mathbb A}
\newcommand{\C}{\mathbb C}
\newcommand{\F}{\mathbb F}
\newcommand{\Gm}{\mathbb G_{\mathrm m}}
\newcommand{\Lef}{\mathbb L}
\newcommand{\PP}{\mathbb P}
\newcommand{\Cl}{\operatorname{Cl}}
\newcommand{\Disc}{\operatorname{Disc}}
\newcommand{\Res}{\operatorname{Res}}
\newcommand{\Sing}{\operatorname{Sing}}
\newcommand{\Sym}{\operatorname{Sym}}
\newcommand{\Spec}{\operatorname{Spec}}
\newcommand{\ord}{\operatorname{ord}}
\newcommand{\id}{\operatorname{id}}
\newcommand{\cE}{E_{\mathrm c}}
\newcommand{\cchi}{\chi_{\mathrm c}}
\newcommand{\set}[1]{\left\{#1\right\}}
\newcommand{\abs}[1]{\left|#1\right|}
\newcommand{\angles}[1]{\left\langle#1\right\rangle}

\crefname{theorem}{theorem}{theorems}
\Crefname{theorem}{Theorem}{Theorems}
\crefname{proposition}{proposition}{propositions}
\Crefname{proposition}{Proposition}{Propositions}
\crefname{lemma}{lemma}{lemmas}
\Crefname{lemma}{Lemma}{Lemmas}
\crefname{corollary}{corollary}{corollaries}
\Crefname{corollary}{Corollary}{Corollaries}
\crefname{question}{question}{questions}
\Crefname{question}{Question}{Questions}


\title[Normal-form complexity]{Normal-Form Complexity of a Three-Sheeted
Keller Cover}
\author{Nathaniel Monson}
\date{July 22, 2026}

\begin{document}

\begin{abstract}
Let \(K=X+Q+C\) be a degree-at-most-three Keller map in dimension \(n\),
with \(Q,C\) homogeneous of degrees two and three.  If the coordinate span of
\(C\) has dimension \(r\), we give a collision-preserving cubic-homogeneous
suspension in dimension \(n+r+1\).  Applied to the public 11-variable
degree-three descendant of the new Jacobian counterexample, where \(r=7\),
this recovers the 19-variable construction independently announced by Harris
Chan.

We determine the generic Jordan type of its nonlinear Jacobian as
\((18,1)\) and give a two-point certificate that its nilpotent flag moves.
Symmetric doubling gives an explicit Hessian-nilpotent quartic in 38
variables, with generic Hessian type \((35,2,1)\).  A two-parameter family
of algebraic inverse rays shows
\[
\Delta^m(\mathcal Q^{m+1})\ne0
\qquad\text{for every }m\ge1,
\]
and identifies the inverse singularities with intersections with the
nonproperness hypersurface.
For the fixed 19-variable cubic tensor, we construct a 110-variable
square-zero cubic-linear pairing and prove that every vector-Waring
decomposition uses at least 52 cubes.  Finally, an exact 13-term functional
is a constant nonzero quartic obstruction across a natural
20-dimensional equivariant family of rank-six cubic compressions.  The
example-specific statements are computer-assisted and come with exact
rational certificates.
\end{abstract}

\maketitle

\begin{center}
\small\emph{Working computer-assisted manuscript.  The recovered exact
certificates have been rerun.  Independent CAS reproduction and specialist
review remain submission requirements.}
\end{center}

\section{Introduction}
\label{sec:introduction}

The classical reductions of Bass--Connell--Wright and Yagzhev make
cubic-homogeneous maps central to the Jacobian conjecture
\cite{bassConnellWright1982}.  Once a three-dimensional counterexample was
announced \cite{alpoge2026announcement}, it therefore became natural to ask
how economically that counterexample reaches the classical normal forms.
We credit Akhil Mathew as the primary human source of the problem: Mathew
prompted Levent Alp\"oge, Alp\"oge prompted Fable, Fable produced the
example, and Alp\"oge announced it publicly
\cite{ulam2026counterexample}.

Atkins-Turkish publicly supplied a degree-at-most-three map in 11 variables,
with three explicit colliding points \cite{atkinsTurkish2026}.  Its cubic
coordinate space has dimension seven.  The construction in
\cref{sec:suspension} consequently gives a cubic-homogeneous counterexample
in
\[
11+7+1=19
\]
variables.  Harris Chan independently announced this special sparse lift on
July 23, 2026 \cite{chan2026nineteen}; we make no sole-priority claim for the
19-variable endpoint.  The reusable contribution is the rank-sensitive
\(n+r+1\) theorem.

The fixed 19-variable tensor has further structure.  Its nonlinear Jacobian
has generic Jordan type \((18,1)\).  Classical symmetric reduction
\cite{deBondtVanDenEssen2005symmetric} then produces a quartic potential in
38 variables.  A 48-variable explicit Hessian witness was already public
through the work of Santibáñez-Leal, starting from Thompson's 24-variable
cubic map \cite{thompson2026twentyfour,santibanezLeal2026cascade}.  Our
claims are the quantitative dimension 38, the exact Hessian Jordan type
\((35,2,1)\), and the all-order Laplacian formula.

In the cubic-linear direction, we use the pairing formalism of Gorni and
Zampieri \cite{gorniZampieri1997}.  The existence of such pairings and the
square-zero mechanism are not new.  For this particular tensor we prove the
model-relative interval
\[
52\le N_{\mathrm{pair}}\le110.
\]
Neither endpoint is asserted to be sharp for arbitrary cubic-homogeneous or
cubic-linear counterexamples.

The main results are summarized as follows.

\begin{theorem}[Rank-sensitive suspension]
\label{thm:intro-suspension}
Let \(K=X+Q+C\colon\A^n\to\A^n\) be Keller over a characteristic-zero
field, with \(Q\) quadratic homogeneous and \(C\) cubic homogeneous.  Put
\[
r=\dim\operatorname{span}\set{C_1,\ldots,C_n}.
\]
Then \(K\) admits a cubic-homogeneous Keller suspension in dimension
\(n+r+1\).  Every collision of \(K\) lifts explicitly to a collision of the
suspension.
\end{theorem}

\begin{theorem}[The fixed 19-variable tensor]
\label{thm:intro-fixed}
For the tensor \(H\) of \cref{sec:fixed-tensor}:
\begin{enumerate}[label=(\roman*)]
\item \(I+H\) is a noninjective cubic-homogeneous Keller map in 19 variables;
\item the generic Jordan type of \(JH\) is \((18,1)\);
\item the associated 38-variable quartic \(\mathcal Q\) is Hessian nilpotent,
has generic Hessian type \((35,2,1)\), and violates Zhao's eventual
vanishing condition at every order; and
\item the vector-Waring length of \(H\) is at least 52, while \(H\) has a
full-rank square-zero pairing of length 110.
\end{enumerate}
\end{theorem}

\begin{theorem}[Equivariant compression obstruction]
\label{thm:intro-obstruction}
Let \(V\) be the space of quadratic source fields that preserve the weights;
it has dimension \(115\).  The locus in \(V\) killing five specified cubic
output rows contains a natural affine slice of dimension \(20\).  Each point
on that slice has at most six linearly independent cubic coordinate
functions.  The same functional, supported on \(13\) terms, takes the value
\(1\) on its quartic obstruction.  Thus no point lifts through quartic order
in the stated equivariant conjugacy model.  The original \(605\)-shift
statement is a sparse corollary.
\end{theorem}

\Cref{thm:intro-obstruction} is deliberately model-relative.  It does not
prove that 19 variables are globally minimal.

\section{A rank-sensitive cubic suspension}
\label{sec:suspension}

Let \(k\) be a field of characteristic zero and write
\[
K(X)=X+Q(X)+C(X),
\]
where \(Q\) and \(C\) are homogeneous of degrees two and three.  Choose a
basis \(q=(q_1,\ldots,q_r)^T\) of the coordinate span of \(C\), and a
constant matrix \(B\in M_{n\times r}(k)\) such that
\[
C=Bq.
\]
Introduce \(w=(w_1,\ldots,w_r)\) and define
\begin{equation}
\label{eq:stable-map}
\widetilde K(X,w)=
\bigl(X+Q(X)+Bw,\ w-q(X)\bigr).
\end{equation}

\begin{lemma}
\label{lem:stable-identity}
The map \(\widetilde K\) is polynomially stably left--right equivalent to
\(K\).
\end{lemma}

\begin{proof}
Set
\[
S(X,w)=(X,w-q(X)),\qquad
T(U,V)=(U+BV,V).
\]
Both are triangular polynomial automorphisms, and direct substitution gives
\[
\widetilde K=T\circ(K\times\id_{\A^r})\circ S.
\]
\end{proof}

Adjoin a homogenizing coordinate \(t\), and put
\begin{equation}
\label{eq:homogeneous-suspension}
G(X,w,t)=
\bigl(
X+tQ(X)+t^2Bw,\ w-q(X),\ t
\bigr).
\end{equation}
Its nonlinear part is cubic homogeneous.

\begin{proof}[Proof of \cref{thm:intro-suspension}]
The Jacobian matrix in the \(X,w\) coordinates has block form
\[
\begin{pmatrix}
I+tJQ(X)&t^2B\\
-Jq(X)&I
\end{pmatrix}.
\]
The Schur complement gives
\begin{align}
\det JG
&=\det\bigl(I+tJQ(X)+t^2BJq(X)\bigr)\notag\\
&=\det JK(tX)=1.\label{eq:schur}
\end{align}
If \(K(p)=K(p')\), then
\[
\widehat p=(p,q(p),1),\qquad
\widehat p'=(p',q(p'),1)
\]
are distinct and satisfy
\[
G(\widehat p)=(K(p),0,1)=G(\widehat p').
\]

Finally, write \(G=I+H\).  Since \(H\) is cubic homogeneous,
\[
JH(sZ)=s^2JH(Z).
\]
The identity \(\det(I+JH)=1\) therefore implies
\(\det(I+\lambda JH)=1\) as a polynomial identity in \(\lambda\).
All nonconstant characteristic coefficients vanish, so \(JH\) is nilpotent.
\end{proof}

\begin{remark}
For a fixed \(K\), any factorization \(C=B'q'\) through \(s\) scalar cubic
coordinates has \(s\ge r\).  Thus \(r\) auxiliaries are optimal within the
factorized suspension \eqref{eq:stable-map}.  This says nothing about other
stable representatives of \(K\) or other homogenization models.
\end{remark}

\section{The fixed 19-variable tensor}
\label{sec:fixed-tensor}

We now work over \(\mathbb Q\).  In the coordinate order
\[
X=(x,y,z,a,b,c,d,q,s,h,k),
\]
let \(\Phi=(\Phi_1,\ldots,\Phi_{11})\) be the public map
\cite{atkinsTurkish2026}:
\begin{align*}
\Phi_1={}&-ac-adz-3ay^2-2az-cd^2+d^2z-ds+7dy^2\\
         &\quad+sxy+3xyz+4y^2+z,\notag\\
\Phi_2={}&-bc-bdz-3by^2-2bz-3cdx-dq+qxy\\
         &\quad+12xy^2+3xz+y,\notag\\
\Phi_3={}&-hk-hxz+kx^2-3x^2y+2x,\\
\Phi_4={}&a-d^2+2dxy,\qquad
\Phi_5=b+3x^2y,\\
\Phi_6={}&c+xyz+3y^2+2z,\qquad
\Phi_7=d-xy,\\
\Phi_8={}&bz+3cx+q,\\
\Phi_9={}&s+az+cxy-xyz-7y^2+cd-dz,\\
\Phi_{10}={}&h-x^2,\qquad
\Phi_{11}=k+xz.
\end{align*}
Its linear part \(L=J\Phi(0)\) has determinant \(-2\).  Normalize by
\[
K=L^{-1}\Phi=X+Q+C.
\]
The coordinate cubics
\[
q_1=C_1,\ q_2=C_2,\ q_3=C_3,\ q_4=C_4,\
q_5=C_5,\ q_6=C_6,\ q_7=C_9
\]
are linearly independent, and the remaining coordinate cubics vanish.  Thus
\(C=Bq\), where \(B\) includes these seven coordinates into
\(\mathbb Q^{11}\).

With \(w=(w_1,\ldots,w_7)\), define
\begin{equation}
\label{eq:H19}
H(X,w,t)=
\bigl(tQ(X)+t^2Bw,\ -q(X),\ 0\bigr)
\end{equation}
in the coordinate order
\[
Z=(x,y,z,a,b,c,d,q,s,h,k,w_1,\ldots,w_7,t).
\]
Then \(G=I+H\) is the suspension \eqref{eq:homogeneous-suspension}.

For reference, three colliding points of \(\Phi\) are
\begin{align*}
p_1={}&(0,0,-1/4,0,0,1/2,0,0,0,0,0),\\
p_2={}&(1,-3/2,13/2,-9/4,9/2,-10,-3/2,3/4,-153/8,1,-13/2),\\
p_3={}&(-1,3/2,13/2,-9/4,-9/2,-10,-3/2,-3/4,-153/8,1,13/2).
\end{align*}
They all map to \((-1/4,0,\ldots,0)\).  Consequently
\[
(p_j,q(p_j),1),\qquad j=1,2,3,
\]
are three colliding points of \(G\).

\subsection{Generic Jordan type}

Let \(A=JH\), and set
\[
E=xy(c+2z)+d^2z+3dy^2.
\]
Define two polynomial vectors
\begin{align}
v_1={}&2tx\,e_h+2tz\,e_k+(xk+zh)e_{w_1},
\label{eq:v1}\\
v_2={}&-txy\,e_a-t(dz+3y^2)e_s-Ee_{w_3}
       +2Ee_{w_6}+xyz\,e_{w_7}.
\label{eq:v2}
\end{align}

\begin{proposition}
\label{prop:jordan19}
Over \(\mathbb Q(Z)\), the matrix \(A\) has Jordan type \((18,1)\).
\end{proposition}

\begin{proof}
Exact differentiation gives
\[
Av_1=Av_2=0.
\]
The vectors are generically independent, while the all-ones specialization
of \(A\) has rank 17.  Hence the generic kernel has dimension two.

Let \(e_t\) be the last coordinate vector.  Exact iteration gives
\begin{equation}
\label{eq:long-chain}
A^{17}e_t=18t^{15}x^9y^5z^2v_2\ne0.
\end{equation}
By \cref{thm:intro-suspension}, \(A\) is nilpotent.  A nilpotent
19-by-19 matrix with a two-dimensional kernel has two Jordan blocks, and
\eqref{eq:long-chain} forces the larger block to have size at least 18.
Thus the partition is \((18,1)\).
\end{proof}

\begin{remark}
At the all-ones specialization the ranks of
\(A,A^2,\ldots,A^{18}\) are
\[
17,16,\ldots,1,0.
\]
The vectors
\[
e_t,Ae_t,\ldots,A^{17}e_t,v_1
\]
specialize to a basis with determinant \(1952152956156672\).
\end{remark}

\section{A restricted compression obstruction}
\label{sec:compression}

Consider quadratic near-identity source conjugacies.  At cubic order their
action is
\[
C\longmapsto C+[Q,P_2],
\qquad
[Q,P_2]=JQ\,P_2-JP_2\,Q.
\]
The shift
\[
P_2=-d^2e_a
\]
lowers the cubic coordinate-span rank from seven to six.  Its exact conjugate
has quartic remainder
\begin{equation}
\label{eq:D4}
(D_4)_z=d^3z+3d^2y^2,\qquad
(D_4)_a=x^2y^2,\qquad
(D_4)_c=-2d^3z-6d^2y^2.
\end{equation}

For a quartic vector field \(R\), define
\begin{align}
\Lambda_4(R)={}&
[x^2y^2]R_a+4[x^2ac]R_a-\frac{20}{3}[x^2as]R_a
+9[x^2qk]R_a\notag\\
&+3[xchk]R_a-\frac12[xd^2q]R_a-7[xshk]R_a
+\frac43[y^2ah]R_a\notag\\
&-\frac12[ybch]R_a-[ydhk]R_a+\frac83[za^2h]R_a
+25[aqhk]R_a-2[dqhk]R_d.
\label{eq:Lambda4}
\end{align}

\begin{proposition}
\label{prop:compression}
For every cubic vector field \(P_3\),
\[
\Lambda_4([Q,P_3])=0,
\qquad
\Lambda_4(D_4)=1.
\]
Thus \(D_4\notin\operatorname{im}[Q,-]\).

Moreover, among the 605 shifts
\[
P_2=\lambda m e_j,
\]
where \(m\) is a quadratic monomial not involving \(X_j\), the shift
\(-d^2e_a\) is the unique one that lowers the cubic span to six.
\end{proposition}

\begin{proof}
The first statement is finite coefficient linear algebra.  Evaluating
\eqref{eq:Lambda4} on all
\[
11\binom{13}{3}=3146
\]
cubic monomial vector fields gives zero, while direct evaluation on
\eqref{eq:D4} gives one.  The restricted homological matrix on the 13
displayed coefficient positions has rank 12 and augmented rank 13.

For the second statement, a fixed nonzero seven-by-seven minor of the cubic
coefficient matrix is evaluated along each of the 605 directions.  Of these,
583 have no exceptional parameter and 22 give 27 rational roots.  Exact
full-rank evaluation at every root leaves only \(P_2=-d^2e_a\).
\end{proof}

\begin{proposition}
\label{prop:target-cleanup}
Use the quadratic output corrections generated by the four outputs
\(d,q,h,k\), which have degree at most two after the conjugacy.  Their ten
quadratic products create a ten-dimensional quartic image.
The \(z\)- and \(c\)-components of \(D_4\) lie outside this image.  The
\(a\)-component has the unique preimage \(-Y_d^2\), but applying it restores
cubic coordinate-span rank seven.
\end{proposition}

\begin{proof}
This is another exact coefficient-matrix calculation.  The ten quartic
images are independent.  Solving against the three nonzero components in
\eqref{eq:D4} gives the assertions.
\end{proof}

\section{The 38-variable Hessian quartic}
\label{sec:hessian}

Let \(F(x)=x+H(x)\) be the 19-variable map.  On
\(\A^{19}_x\times\A^{19}_y\), define the cotangent double
\[
\Phi_F(x,y)=\bigl(F(x),JF(x)^Ty\bigr).
\]
Its Jacobian is block lower triangular, with diagonal blocks \(JF\) and
\(JF^T\), so \(\det J\Phi_F=1\).

Over \(\mathbb Q(i)\), put
\[
T(u,v)=\left(u+iv,\frac{u-iv}{2}\right).
\]
The scalar block satisfies
\[
TT^T=
\begin{pmatrix}0&1\\1&0\end{pmatrix}.
\]
Therefore
\[
\mathcal G=T^{-1}\Phi_FT=I+\nabla\mathcal Q,
\]
where
\begin{equation}
\label{eq:Q38}
\mathcal Q(u,v)=\frac12(u-iv)^TH(u+iv).
\end{equation}
This is a homogeneous quartic with 724 nonzero monomials.  If
\(F(p)=F(p')\), then
\[
\Phi_F(p,0)=\Phi_F(p',0),
\]
so \(\mathcal G\) has an explicit transported collision.

\begin{proposition}
\label{prop:hessian-jordan}
The generic Jordan type of \(\operatorname{Hess}\mathcal Q\) is
\[
(35,2,1).
\]
\end{proposition}

\begin{proof}
Before the linear twist, the nonlinear Jacobian is
\[
N=
\begin{pmatrix}
A&0\\
\mathcal L&A^T
\end{pmatrix},
\qquad
A=JH(x),\qquad
\mathcal L=\operatorname{Hess}_x(y^TH(x)).
\]
The Hessian of \(\mathcal Q\) is similar to \(N\).

The kernel vectors \eqref{eq:v1}--\eqref{eq:v2} satisfy
\[
D^2H[v_1,v_2]=D^2H[v_2,v_2]=0.
\]
Since they span \(\ker A\) generically, this implies
\(\mathcal Lv_2\in\operatorname{im}A^T\).  In addition to the two kernel
vectors \((0,\ker A^T)\), the matrix \(N\) therefore has a kernel vector
whose first component is \(v_2\).  Hence \(\operatorname{rank}N\le35\).
An exact all-ones specialization has rank 35.

For \(k\ge1\), the lower-left block of \(N^k\) is
\[
\sum_{j=0}^{k-1}(A^T)^j\mathcal L A^{k-1-j}.
\]
Since \(A^{18}=0\), the only possible term for \(k=35\) is
\((A^T)^{17}\mathcal L A^{17}\).  The image of \(A^{17}\) is spanned by
\(v_2\), and \(D^2H[v_2,v_2]=0\), so this term vanishes.  Thus \(N^{35}=0\).
At the all-ones specialization,
\[
(N^{34})_{38,19}=648\ne0.
\]
The nilpotency index is 35 and the rank is 35, so there are three blocks.
The remaining three dimensions split as \(2+1\), proving the result.
\end{proof}

\subsection{An all-order Laplacian certificate}

Let
\[
\bar Y=(1/2,1,0,\ldots,0,1)\in\mathbb Q^{19},
\qquad
\bar W=T^{-1}(\bar Y,0)
=\left(\frac{\bar Y}{2},-\frac{i\bar Y}{2}\right).
\]
An exact inverse branch of the original three-dimensional map, lifted through
the stable reductions, gives
\begin{equation}
\label{eq:inverse-ray-G}
\bigl(F^{-1}(s\bar Y)\bigr)_1
=\frac{1/\sqrt{1-s^4}-1}{s^3}
=\sum_{r\ge1}\frac1{4^r}\binom{2r}{r}s^{4r-3}.
\end{equation}
Consequently
\begin{equation}
\label{eq:inverse-ray-Q}
\bigl(\mathcal G^{-1}(s\bar W)\bigr)_{u_1}
=\frac12\sum_{r\ge1}\frac1{4^r}\binom{2r}{r}s^{4r-3}.
\end{equation}

Put \(P=-\mathcal Q\), so that
\(\mathcal G=I-\nabla P\).  Zhao's inversion formula for a
Hessian-nilpotent potential gives homogeneous inversion terms
\cite{zhao2007hessian}
\[
R_m=
\frac{1}{2^m m!(m+1)!}\Delta^m(P^{m+1}).
\]
The degree of \(\nabla R_m\) is \(2m+3\).  Comparing
\eqref{eq:inverse-ray-Q} at \(m=2r-3\) gives:

\begin{theorem}
\label{thm:laplacian}
For every \(r\ge2\),
\[
\partial_{u_1}\Delta^{2r-3}
\bigl(\mathcal Q^{2r-2}\bigr)(\bar W)
=
\frac{(2r-3)!(2r-2)!}{16}\binom{2r}{r}\ne0.
\]
In particular,
\[
\partial_{u_1}\Delta(\mathcal Q^2)(\bar W)=\frac34.
\]
\end{theorem}

\begin{proof}
The exponent \(m+1=2r-2\) is even, so
\(P^{m+1}=\mathcal Q^{m+1}\).  Multiply the coefficient of
\(s^{4r-3}\) in \eqref{eq:inverse-ray-Q} by
\(2^m m!(m+1)!\).
\end{proof}

Since \(\mathcal Q\) is Hessian nilpotent,
\(\Delta^m(\mathcal Q^m)=0\) for all \(m\ge1\), while
\cref{thm:laplacian} gives infinitely many nonzero
\(\Delta^m(\mathcal Q^{m+1})\).  This is an explicit failure of eventual
vanishing in Zhao's framework.

\section{A \texorpdfstring{\(52\)-to-\(110\)}{52-to-110}
square-zero pairing interval}
\label{sec:pairing}

Let \(V=\mathbb Q^{19}\), and regard
\[
H\in V\otimes\Sym^3(V^*).
\]
Its \emph{vector-Waring length} is the least \(R\) for which
\begin{equation}
\label{eq:waring}
H(z)=\sum_{\nu=1}^R b_\nu\ell_\nu(z)^3.
\end{equation}

\subsection{The lower bound}

Let
\[
W=\operatorname{span}\set{
\partial H_i/\partial z_j
}\subset\Sym^2(V^*).
\]
Every decomposition \eqref{eq:waring} gives
\[
W\subset U:=\operatorname{span}\set{\ell_\nu^2}.
\]
Exact row reduction gives
\begin{equation}
\label{eq:dimW}
\dim W=39.
\end{equation}

In the coordinate order of \cref{sec:fixed-tensor}, the 12-dimensional
subspace
\[
\mathcal A=
\langle z,c,q,s,k,w_1,\ldots,w_7\rangle
\]
is a common zero block:
\[
W|_{\mathcal A}=0.
\]
The coefficient system for a constant vector \(v\) satisfying
\(JH(z)v=0\) has rank 19; a specified 19-by-19 minor has determinant 18.
Thus the \(\ell_\nu\) span \(V^*\), and their restrictions span
\(\mathcal A^*\).  Twelve independent restrictions have twelve independent
squares.  It follows that
\[
\dim(U/W)\ge12,\qquad \dim U\ge51.
\]

\begin{proposition}
\label{prop:waring52}
Every vector-Waring decomposition of \(H\) has length at least 52.
\end{proposition}

\begin{proof}
Suppose a decomposition has length at most 51.  Every preceding inequality is
then an equality:
\[
R=\dim U=51,\qquad W=\ker\bigl(U\to\Sym^2(\mathcal A^*)\bigr).
\]
Choose a basis
\(\alpha_1,\ldots,\alpha_{12}\) of \(\mathcal A^*\) among the restricted
forms.  If
\(\alpha=\sum c_j\alpha_j\) and
\(\alpha^2\) lies in the span of the \(\alpha_j^2\), then all
\(c_ic_j\) for \(i\ne j\) vanish.  Hence every restricted form is
proportional to one \(\alpha_j\).

Choose a complement \(V=\mathcal B\oplus\mathcal A\), with
\(\dim\mathcal B=7\), and let
\[
\mathcal C=
\operatorname{pr}_{\mathcal B^*\otimes\mathcal A^*}(W).
\]
The equality description forces \(\mathcal C\) to be invariant under the
twelve coordinate projectors in the \(\alpha_j\)-basis.  Therefore its right
stabilizer algebra
\[
\mathcal E=
\set{T\in\operatorname{End}(\mathcal A^*):
(I\otimes T)\mathcal C\subset\mathcal C}
\]
contains a conjugate of the full diagonal algebra.  Its commutant can then
contain no nonzero nilpotent.

Exact rational computation gives
\[
\dim\mathcal C=24,\qquad
\dim\mathcal E=62,\qquad
\dim\mathcal E'=7,
\]
and \(\mathcal E'\) contains the matrix unit \(E_{5,0}\), with indices
numbered \(0,\ldots,11\).  This element is nonzero and square-zero, a
contradiction.
\end{proof}

\subsection{The 110-variable upper construction}

The supplement contains rational matrices
\[
B\in M_{19\times110}(\mathbb Q),\qquad
D\in M_{110\times19}(\mathbb Q)
\]
satisfying
\begin{equation}
\label{eq:pair-data}
\operatorname{rank}B=\operatorname{rank}D=19,\qquad
BD=0,\qquad
H(z)=B(Dz)^{*3}.
\end{equation}
Put \(A=DB\) and
\[
F_A(W)=W+(AW)^{*3}.
\]
Then \(A^2=0\) and \(\operatorname{rank}A=19\).

\begin{proposition}
\label{prop:pair110}
The map \(F_A\colon\A^{110}\to\A^{110}\) is a noninjective Keller map.
It is Gorni--Zampieri paired with \(G=I+H\).
\end{proposition}

\begin{proof}
Choose a right inverse \(C\) of \(B\).  The exact data give \(AC=D\), and
\[
B F_A(Cz)=z+B(Dz)^{*3}=G(z).
\]
By Sylvester's determinant identity,
\begin{align*}
\det JF_A(W)
&=\det\left(I_{110}
+3\operatorname{diag}((AW)^2)DB\right)\\
&=\det\left(I_{19}
+3B\operatorname{diag}((DBW)^2)D\right)\\
&=\det JG(BW)=1.
\end{align*}

If \(G(z_1)=G(z_2)\), put \(u=Cz_1\), \(v=Cz_2\), and
\(\delta=F_A(u)-F_A(v)\).  Then \(B\delta=0\), hence
\(A\delta=DB\delta=0\).  Thus
\[
F_A(v+\delta)=F_A(v)+\delta=F_A(u).
\]
The exact collision certificate verifies \(u\ne v+\delta\).
\end{proof}

\begin{corollary}
\label{cor:pair-interval}
For the fixed tensor \(H\), the minimum length of a full-rank square-zero
pairing satisfies
\[
52\le N_{\mathrm{pair}}\le110.
\]
\end{corollary}

\begin{proof}
The lower bound is \cref{prop:waring52}; every such pairing is a
vector-Waring decomposition.  The upper bound is \cref{prop:pair110}.
\end{proof}

\section{Exact computation and reproducibility}
\label{sec:computation}

All example-specific calculations use exact rational arithmetic.  The fresh
audit reran:
\begin{enumerate}[label=(\arabic*)]
\item the public-input-to-19D collision and nilpotency certificate;
\item the rank-seven factorization and Schur-complement identities;
\item both kernel vectors, the length-18 Jordan chain, and the degeneration
certificate;
\item the 38D gradient identity, Hessian rank and index, inverse branch, and
first Laplacian value;
\item the 110D matrices, square-zero identity, pairing identity, Keller
determinant transfer, and collision;
\item the 52-cube lower-bound matrices and equality obstruction;
\item all 3,146 monomial tests for \(\Lambda_4\);
\item all 605 triangular monomial shifts and the ten target-cleanup
directions; and
\item a separate verifier using only Python rational fractions for the small
rank, minor, circuit, and commutant certificates.
\end{enumerate}

The regenerated 724-term quartic is byte-for-byte identical to the recovered
output.  The 110D pairing JSON and the serialized small certificates also
match their archived inputs exactly.

The source, exact serialized tensors, and successful logs are organized in
\texttt{computational-supplement/}.  In particular,
\texttt{fixed-tensor/} contains the 19D, 38D, and 110D constructions,
\texttt{continuation/} contains the Waring and compression calculations, and
the two \texttt{n5-regular-*} directories separate the collision-line
existence theorem from the first-normal global obstruction.  The
claim-to-file map and minimal replay commands are in
\texttt{COMPUTATION.md}.  This organization reflects the mathematical
lesson of the calculation: regularity on the collision line is possible,
but global nilpotence imposes new normal-direction equations.

The computations do not establish novelty, and the several SymPy scripts are
not independent CAS implementations.  The Fraction-only checker is
independent of SymPy but consumes serialized certificate matrices.  A public
release should additionally reproduce the main identities in Singular,
Macaulay2, Sage, or Magma.

\section{Questions}
\label{sec:questions}

\begin{question}[Minimum homogeneous dimension]
\label{q:min-homogeneous-dimension}
What is the minimum dimension of a cubic-homogeneous counterexample?  The
fixed \(19\)-variable construction gives an upper bound, while known
low-dimensional classifications and the five-dimensional restrictions
collected in the companion research register leave a substantial gap.
\end{question}

\begin{question}
Can one improve the rank-sensitive complexity
\[
\dim K+\dim\operatorname{span}\set{(K_3)_i}
\]
by choosing a different degree-at-most-three stable representative?
\end{question}

\begin{question}[Intrinsic tensor and pairing ranks]
\label{q:intrinsic-ranks}
What are the exact vector-Waring length, pairing rank, and square-zero
realization rank of the fixed \(19\)-variable cubic tensor?  In particular,
can the present vector-Waring interval between \(52\) and \(110\) be closed
without committing to one reduction algorithm?
\end{question}

\begin{question}
Can the \(\mathfrak{sl}_{18}\) collision-monolith structure recorded in the
companion research register be used to derive the Jordan types
\((18,1)\) and \((35,2,1)\) without coefficient calculation?
\end{question}

\appendix
\input{appendices/structural-extensions}
\section{Full-kernel pencils and the \texorpdfstring{\(19\to18\)}{19-to-18}
  compression problem}
\label{app:full-kernel-and-compression}

Two computations continue our dimension-five and nineteen-variable analyses.
The first reorganizes the regular full-kernel pencil by binary-sextic
invariant theory.  The second proves a no-lift theorem on a natural
transversal to the \(19\to18\) compression locus.

\subsection{The regular full-kernel chart}

In the full-kernel chart write
\[
M(s,t)=R(s,t)T^{-1}K(s,t).
\]
With the standard quartic Veronese column \(k\), cubic column \(q\), and
row \(\ell\), exact multiplication gives
\[
N_4^3=6q\ell,\qquad Rq=a_7k,
\]
and hence
\[
M(s,t)^4=6a_7\,k(s,t)\ell(s,t)T^{-1}K(s,t).
\]

\begin{proposition}[Exact regularity criterion]
\label{prop:full-kernel-regularity}
On \(\det T\ne0\), the pencil \(M(s,t)\) has Jordan type \((5)\) at every
point of \(\PP^1\) if and only if \(a_7\ne0\).
\end{proposition}

\begin{proof}
All factors in the displayed formula other than \(a_7\) are nonzero at
every point of \(\PP^1\), while \(M^5=0\).  Thus \(M^4\) is everywhere
nonzero exactly when \(a_7\ne0\).
\end{proof}

Put \(\lambda=a_7\).  After a fixed linear change of the other seven
parameters,
\[
T=\frac{\lambda}{4}I+S_b,
\qquad
S_b^TJ+JS_b=0,
\]
and the \(b\)-space is the irreducible principal-\(\mathfrak{sl}_2\) module
\[
\Sym^6(\C^2).
\]
Thus the chart is naturally
\[
\C\lambda\oplus\Sym^6(\C^2),
\]
not an unstructured eight-parameter matrix family.  If \(A,B\) are the
standard degree-two and degree-four binary-sextic invariants in the chosen
normalization, then
\[
\det T=\frac{(\lambda^2+4A)^2-96B}{256}.
\]

The marking equations are linear in the nine unfixed marking variables.
On an explicit pivot chart they eliminate to one degree-\(12\) parameter
equation and five degree-\(13\) equations.  A candidate invariant formula
for the degree-\(12\) equation was reconstructed by exact modular
interpolation and checked on rational samples; it is retained in the
supplement as a computational conjecture until a coefficient certificate or
conceptual invariant-theoretic derivation is supplied.

\begin{proposition}[Finite-field sample audit]
\label{prop:full-kernel-finite-field-samples}
An exhaustive \(\mathbf F_{13}\)-calculation finds thirteen rational points
of this normalized collision system.  Together with the previously checked
six points over \(\mathbf F_7\) and eleven over \(\mathbf F_{11}\), all
thirty sampled points have nonzero first-normal obstruction.  This is exact
finite-field evidence, not a proof that the obstruction section is nowhere
zero on the characteristic-zero collision curve.
\end{proposition}

\begin{proof}[Exact computation]
The finite-field programs enumerate the normalized solutions and evaluate the
first-normal obstruction in the corresponding prime field.  Their saved exact
outputs give the three counts and a nonzero value at every listed point.
\end{proof}

\subsection{The unrestricted compression tangent space}

Let \(K=X+Q+C\colon\A^{11}\to\A^{11}\) be the normalized
degree-at-most-three map used in the fixed nineteen-variable construction.
A quadratic source field \(P\) changes its cubic term to
\[
C_P=C+[Q,P].
\]

\begin{proposition}[Large row-killing families]
\label{prop:row-killing-families}
The affine space of quadratic source fields has dimension \(726\).
There is a \(109\)-dimensional affine family on which five prescribed cubic
output rows vanish and the remaining coordinate span has rank at most six.
The quartic obstruction functional is identically one on this family.
Inside it, a specified \(75\)-dimensional family consists of triangular
polynomial automorphisms.  The full rank-six tangent space at the base point
has dimension \(129\), leaving a \(20\)-dimensional quotient by the
row-killing directions; a nonempty Zariski-open subset of that quotient is
obstructed at second order.
\end{proposition}

\begin{proof}[Exact computer-assisted proof]
The assertions are rational ranks and symbolic bracket identities in the
displayed coefficient bases.  Fixed full-rank minors certify the ranks, and
the triangular family is obtained by removing the \(145\) forbidden output
coefficients from the \(220\)-dimensional triangular ansatz.
\end{proof}

\subsection{The complete transverse cone}

At the row-killing rank-six point
\[
P_*=2hx\,e_z-d^2e_a-6ax\,e_b+
(2az-\tfrac23bh)e_c+2b^2e_s-h^2e_k,
\]
the determinantal tangent space has dimension \(115\).  Modulo the
\(109\) row-killing directions, choose transverse coordinates
\(u_0,\ldots,u_5\).

\begin{theorem}[Transverse no-lift]
\label{thm:transverse-cone-no-lift}
In this six-dimensional transversal, the second-order rank-six cone is the
reduced scheme
\[
V(u_0u_5,u_2u_3,u_3u_5)
=
V(u_3,u_5)\cup V(u_0,u_3)\cup V(u_2,u_5).
\]
There is an explicit polynomial family \(P(u)\) such that
\[
\operatorname{rank}(C+[Q,P(u)])\le6
\quad\Longleftrightarrow\quad
u_0u_5=u_2u_3=u_3u_5=0.
\]
Every point of the cone is therefore integrated to an exact cubic jet of
coordinate span six.  Nevertheless no member lifts through quartic order.
\end{theorem}

\begin{proof}[Exact computer-assisted proof]
The projected second fundamental form has the three displayed generators.
Five exact pivot charts show that the Schur-complement ideal of the explicit
family is exactly their radical intersection and that six output rows retain
rank six everywhere on it.

For a quartic vector field \(E\), set
\[
\Gamma_4(E)=[h^4](E_z).
\]
Support alone gives \(\Gamma_4([Q,P_3])=0\) for every cubic source
correction \(P_3\).  Direct exact calculation on the whole six-parameter
family gives
\[
\Gamma_4(\mathcal O_4(P(u)))=-1.
\]
Thus the quartic homological equation cannot hold at any rank-six member.
\end{proof}

The theorem resolves the complete natural six-dimensional transversal, but
not the entire \(726\)-parameter compression problem: nonlinear coupling
with the \(109\) row-killing parameters, independent target changes, and
stable-coordinate changes remain outside its scope.  A previously studied
two-parameter active surface is superseded by
\cref{thm:transverse-cone-no-lift} and is not used.

\section*{Companion results and research register}

This reader edition is accompanied by a separate \emph{Results and Research
Register}.  The register preserves secondary proved results, conditional
developments, open questions, corrections, and computational leads without
making them part of this paper's theorem spine.  Proof-bearing technical
developments omitted from this reader PDF remain in the versioned source
tree and in the complete version-8 archival edition.  The v9 primary-home
map distinguishes the paper responsible for a result from other papers that
use or motivate it.


\section*{AI and computational disclosure}

AI systems were used extensively in exploration, symbolic-program
development, proof drafting, source organization, and manuscript editing.
They are not authors.  The public 11-variable source also discloses that its
explicit construction and verification code were generated by ChatGPT
\cite{atkinsTurkish2026}.  The evidence for the computer-assisted theorems is
the exact certificate suite described in \cref{sec:computation}, not language
model output.  Human specialist review has not yet been completed.

\bibliographystyle{amsplain}
\bibliography{../common/references}

\end{document}
